\documentclass[11pt]{article}
\usepackage[T1]{fontenc}
\usepackage{lmodern}
\usepackage[margin=1in]{geometry}
\usepackage{amsmath,amssymb,amsthm,mathtools}
\usepackage[hidelinks]{hyperref}
\usepackage{microtype}
\newtheorem{lemma}{Lemma}[section]
\newtheorem{corollary}[lemma]{Corollary}
\setlength{\parindent}{0pt}
\setlength{\parskip}{5pt}
\emergencystretch=2em
\title{Targeted Lean 4 blueprint\\\large Repairs to the three blocked spots at lines 2084--2086}
\author{}
\date{}
\begin{document}
\maketitle
This document concerns only the blocked proof in
\path{task_2_smoothingineq3d_blueprint_updated.tex}.
Declaration names below are proposed project names, not claims about
existing Mathlib identifiers. This is a mathematical blueprint; no Lean
compiler has checked it.

Sections~1 and~2 give independently usable local repairs. For the merged
repair, use Section~3 instead of the entire old stopping-time proof: its
measurable-set refinements do not require either spatial dyadic resolution
or principal cubes. It also replaces the overly general lemma statement
by the actual monomial instance. The supplied replacement fragment covers
the old lemma and its immediate corollary (original lines 2069--2136), and
leaves the rest of the blueprint unchanged.

\section{Spot 1: a finite dyadic tree cannot resolve arbitrary simple inputs}
\textbf{Blocked sentence.} ``Inside each enlargement take the finite dyadic
tree resolving all level sets of the three nonnegative simple inputs.''

For a local repair retaining a spatial-tree argument, first take
nonnegative bounded inputs supported in a bounded half-open root cube $R$.
For its depth-$n$ dyadic partition $\mathcal D_n(R)$ define
\[
 f_{i,n}=\sum_{Q\in\mathcal D_n(R)}
       \left(|Q|^{-1}\int_Q f_i\right)\mathbf1_Q,
\]
extended by zero outside $R$. The finite tree of generations $0,\ldots,n$
resolves the level sets of $f_{i,n}$, not of $f_i$. The approximation depth
$n$ must remain an explicit parameter.

\paragraph{Lean proof obligations.}
\texttt{dyadicApprox\_basic}: prove finite range, measurability,
nonnegativity, support in $R$, and $\|f_{i,n}\|_\infty\leq\|f_i\|_\infty$.

\texttt{dyadicApprox\_tendsto\_L1}: prove
$\|f_{i,n}-f_i\|_1\to0$. Averaging is an $L^1$-contraction. Approximate
$f_i$ in $L^1(R)$ by a continuous function on $\overline R$, and use uniform
continuity to obtain uniform convergence of its dyadic averages. For
$1\leq p<\infty$, boundedness also gives
\[
 \|f_{i,n}-f_i\|_p^p\leq
 (2\|f_i\|_\infty)^{p-1}\|f_{i,n}-f_i\|_1\longrightarrow0.
\]
Handle the zero bound separately.

\texttt{translationAverage\_tendsto\_L1\_of\_boundedApprox}: for a normalized
positive translation average, the telescoping expansion and the trivial
$L^1\times L^\infty\times L^\infty\to L^1$ bounds give
\[
 \|T_1(f_{1,n},f_{2,n},f_{3,n})-T_1(f_1,f_2,f_3)\|_1
 \leq\sum_i\|f_{i,n}-f_i\|_1\prod_{a\ne i}\|f_a\|_\infty.
\]
Hence the outputs converge in $L^1$ and have an almost-everywhere
convergent subsequence, independently of any subunit bound.

\texttt{dyadic\_bound\_passes\_to\_bounded\_inputs}: if the finite-tree bound
has a constant independent of $n$, apply Fatou to the $q$-th powers along
that subsequence, and use the input $L^{q_i}$ convergence. The dyadic
averages need not be monotone. Only the subsequent extension using
\[
 f_i^{(K)}=\mathbf1_{[-K,K)^3}\min(f_i,K)
\]
is monotone, and permits monotone convergence in the inputs, positive
average, and $q$-th power.

\section{Spot 2: principal cubes, averages, and the correct packing quantity}
\textbf{Blocked sentence.} ``Call $Q$ principal ... if its local
$L^{r_i}$-mass is more than twice the sum of the masses of its strict
descendants.''

Fix a finite dyadic tree. Explicitly supply exponents $s_i>1$ for the
three original inputs and distinguish them from the adjoint exponents
indexed by $0,j',k$. Define
\[
 M_i(Q)=\int_Q f_i^{s_i},\qquad a_i(Q)=M_i(Q)/|Q|.
\]
Make the root principal. For a principal $P$, take as stopping children
the maximal proper descendants satisfying
$\exists i,\ a_i(Q)>6a_i(P)$. Recursively declare those children principal.
For a selected child set
\[
 \operatorname{colour}(P,Q)=\min\{i:a_i(Q)>6a_i(P)\}.
\]
Here descendants are searched in the specified finite tree; sums are over
selected stopping children, not all descendants.

\paragraph{Lean proof obligations.}
\texttt{stopChildren\_pairwise\_disjoint}: use dyadic nested-or-disjoint
geometry and maximality. \texttt{principal\_family\_laminar}: prove only
nested-or-disjoint, not that the entire family is a chain.

\texttt{stopChildren\_volume\_le\_half}: for each colour $i$, disjointness
and the defining threshold give
\[
 6a_i(P)\sum_{\operatorname{colour}(P,Q)=i}|Q|
 \leq\sum_{\operatorname{colour}(P,Q)=i}M_i(Q)\leq M_i(P).
\]
If $a_i(P)>0$, divide to obtain $|P|/6$. If $a_i(P)=0$, there are no
children of that colour. Summing the three colour bounds gives
\[
 \sum_{Q\in\operatorname{Stop}(P)}|Q|\leq |P|/2.
\]

\texttt{principal\_good\_regions}: define
$E_P=P\setminus\bigcup_{Q\in\operatorname{Stop}(P)}Q$.
Prove $|E_P|\geq|P|/2$, pairwise disjointness of the $E_P$, and
$\sum_{P\subseteq R}|P|\leq2|R|$ for principal cubes. If the inputs are
resolved at the terminal generation, then
$f_i^{s_i}\leq6a_i(P)$ almost everywhere on $E_P$; otherwise a terminal
cube would qualify and lie under a maximal qualifying descendant.

Recursion terminates by decreasing remaining dyadic depth. Comparisons
and finite maximal selection may be classical/noncomputable.
\textbf{Do not infer a mass contraction.} All of $M_i(P)$ may be carried
by one small stopping child. The proved contraction is for Lebesgue
volume, and does not establish the old $2^{-m}$ analytic remainder.

\section{Spot 3: replace the unspecified cycles by an explicit monomial proof}
\textbf{Blocked sentence.} ``Iterate this operation cyclically through the
three inputs. After $m$ cycles, expansion by positivity gives ...''

Delete the asserted absorption inequality and the formulas for $d_m$ and
$q_{i,m}$. They do not follow from the local repairs in Sections~1--2.
Moreover, the abstract adjoint-gain implication as stated is false; a
verification is provided at the end. The replacement below proves the
actual monomial endpoint with $q=8/9$ and preserves the distinguished
$L^2$ input.

The source PDF, p.~79, states (6.24) and attributes it to Corollary~5.3,
but does not provide the proposed cube recursion or exponent formulas.
The replacement instead adapts the explicit measurable refinements of
Section~5.3 and the parameter-tower/Jacobian mechanism of Sections~5.4
and~5.6. It is a new three-set specialization of that mechanism, fully
specified here.

% Targeted replacement for the old lemma and immediate corollary,
% original source lines 2069--2136. The historical lemma label is retained
% for reference compatibility; its statement is deliberately specialized.
% This is a mathematical Lean 4 blueprint, not compiler-checked Lean code.

\begin{lemma}[Three-set restricted estimate for the upper-half monomial average]
\label{lem:three-set-restricted-upper-half}
Let $I=(1/2,1)$, $\gamma_i(t)=t^i e_i$ for $i\in\{1,2,3\}$, and
\[
 T(f_1,f_2,f_3)(x)=2\int_I\prod_{i=1}^3 f_i(x+\gamma_i(t))\,dt.
\]
For $d\in\{2,3\}$ put $\ell=5-d$. For measurable sets of finite measure,
\[
 K:=\int_{\mathbb R^3}T(\mathbf1_{E_1},\mathbf1_{E_2},\mathbf1_{E_3})(x)\,dx
 \leq 16|E_1|^{1/2}|E_d|^{1/2}|E_\ell|^{1/4}.
 \tag{R}
\]
\end{lemma}

\begin{proof}
This is a three-set specialization of the measurable-refinement and real
parameter-flow mechanism in Kosz et al., Section~5.3, Lemma~5.14, and
Section~5.6. The specific flow and constants below are derived here; they
are not a transcription of the assertion surrounding equation~(6.24).

\paragraph{Incidence and the exact deletion identity.}
Work first with Borel representatives. Changes on null sets do not change
any incidence integral, by translation invariance and Tonelli. For a tuple
$F=(F_1,F_2,F_3)$ define
\[
 \mathcal K(F)=2\int_I\int_{\mathbb R^3}
       \prod_{i=1}^3\mathbf1_{F_i}(x+\gamma_i(t))\,dx\,dt,
\]
\[
 H_i[F](y)=2\int_I\prod_{a\ne i}
 \mathbf1_{F_a}\bigl(y-\gamma_i(t)+\gamma_a(t)\bigr)\,dt.
\]
The functions $H_i[F]$ are measurable and lie in $[0,1]$. Translation in
$x$ gives the identity
\[
 \mathcal K(F)=\int_{F_i}H_i[F](y)\,dy.                 \tag{D1}
\]
These $H_i$ are two-factor fibre-incidence functions, not the three-input
adjoints of the original blueprint.

If $K=0$, (R) is immediate. Otherwise $0<|E_i|<\infty$ for every $i$ and
\[
 \alpha_i=K/|E_i|\in(0,1].
\]
For $F_i\subseteq E_i$ and $c>0$, replace only $F_i$ by
\[
 F_i'=\{y\in F_i:H_i[F](y)\geq c\alpha_i\}.
\]
Keep the other two sets unchanged. Equation (D1) proves
\[
 \mathcal K(F)-\mathcal K(F')
 =\int_{F_i\setminus F_i'}H_i[F]
 \leq c\alpha_i|E_i|=cK.                             \tag{D2}
\]
In Lean, prove (D1) as \texttt{incidence\_eq\_fibreIntegral} and (D2) as
\texttt{refine\_one\_loss}. The reference mass $K$ and the $\alpha_i$
remain those of the original sets at every stage.

\paragraph{Exactly three updates.}
Let $F^{(0)}=(E_1,E_2,E_3)$. Define successively
\[
 F^{(1)}_1=\{y\in E_1:H_1[F^{(0)}](y)\geq\alpha_1/2\},
 \qquad F^{(1)}_a=E_a\quad(a\ne1),
\]
\[
 F^{(2)}_d=\{y\in F^{(1)}_d:H_d[F^{(1)}](y)\geq\alpha_d/4\},
 \qquad F^{(2)}_a=F^{(1)}_a\quad(a\ne d),
\]
\[
 F^{(3)}_1=\{y\in F^{(2)}_1:H_1[F^{(2)}](y)\geq\alpha_1/8\},
 \qquad F^{(3)}_a=F^{(2)}_a\quad(a\ne1).
\]
Thus the update order is $1,d,1$, with thresholds $1/2,1/4,1/8$.
Three applications of (D2) give
\[
 \mathcal K(F^{(1)})\geq K/2,\qquad
 \mathcal K(F^{(2)})\geq K/4,\qquad
 \mathcal K(F^{(3)})\geq K/8.                        \tag{D3}
\]
In particular $F^{(3)}_1$ is nonempty. Choose $z\in F^{(3)}_1$ once.
This is a finite choice of a point, not measurable selection of a family.

\paragraph{Explicit tower and routing.}
Define
\[
 U=\{u\in I:\ \forall a\ne1,\
       z-\gamma_1(u)+\gamma_a(u)\in F^{(2)}_a\}.
\]
For every real $u$ set $y_1(u)=z-\gamma_1(u)+\gamma_d(u)$ and define
\[
 V(u)=\{v\in I:\ \forall a\ne d,\
       y_1(u)-\gamma_d(v)+\gamma_a(v)\in F^{(1)}_a\}.
\]
For every real pair $(u,v)$ set $y_2(u,v)=y_1(u)-\gamma_d(v)+\gamma_1(v)$ and define
\[
 W(u,v)=\{w\in I:\ \forall a\ne1,\
       y_2(u,v)-\gamma_1(w)+\gamma_a(w)\in E_a\}.
\]
The three threshold definitions, read in reverse chronological order,
prove
\[
 |U|\geq\alpha_1/16,\qquad
 |V(u)|\geq\alpha_d/8,\qquad
 |W(u,v)|\geq\alpha_1/4.                            \tag{T1}
\]
The $V$ bound is for $u\in U$, and the $W$ bound is for $u\in U$, $v\in V(u)$. The factor $2$ in the definition of $H_i$ accounts for these denominators.
The memberships used are precisely
\[
 z\in F^{(3)}_1
 \longmapsto y_1\in F^{(2)}_d
 \longmapsto y_2\in F^{(1)}_1
 \longmapsto \Phi(u,v,w)\in E_\ell.
\]
Remove a single interval from the middle fibre:
\[
 V'(u)=V(u)\cap\{v:|u-v|\geq\alpha_d/32\}.
\]
The removed interval has length $\alpha_d/16$, so $|V'(u)|\geq\alpha_d/16$.
Now set
\[
 \mathcal P=\{(u,v,w)\in I^3:\ u\in U,\ v\in V'(u),\ w\in W(u,v)\}.
\]
All sets are given by explicit measurable predicates. Tonelli and (T1)
yield
\[
 |\mathcal P|\geq 2^{-10}\alpha_1^2\alpha_d.         \tag{T2}
\]
Prove this as \texttt{three\_step\_tower}. No stopping cubes, height
induction, or unidentified coordinate choices occur.

\paragraph{The polynomial flow.}
Put
\[
 \Phi(u,v,w)=z-\gamma_1(u)+\gamma_d(u)
                -\gamma_d(v)+\gamma_1(v)
                -\gamma_1(w)+\gamma_\ell(w).
\]
The routing above proves $\Phi(\mathcal P)\subseteq E_\ell$. In coordinate
order $(1,d,\ell)$,
\[
 \Phi(u,v,w)
 =\bigl(z_1-u+v-w,\ z_d+u^d-v^d,\ z_\ell+w^\ell\bigr).
\]
Its absolute Jacobian is
\[
 |\det D\Phi|=
 \begin{cases}
  6w^2|u-v|,&d=2,\\
  6w(u+v)|u-v|,&d=3.
 \end{cases}                                                    \tag{J1}
\]
Since $u,v,w\in I$, both expressions are at least $|u-v|$ and consequently
at least $\alpha_d/32$ on $\mathcal P$.

The map is injective on $\Omega=\{(u,v,w)\in I^3:u\ne v\}$. Indeed its
$\ell$-coordinate determines $w>0$. Its first coordinate then determines
$h=u-v\ne0$. If $d=2$, its $d$-coordinate minus $z_d$, divided by $h$, determines $u+v$.
If $d=3$, it determines $c=(u^3-v^3)/(u-v)$, and
\[
 4c=3(u+v)^2+h^2
\]
determines $u+v>0$. In either case $u$ and $v$ follow from their sum and
difference. This avoids a general polynomial fibre-counting theorem.

Prove \texttt{three\_step\_flow\_injective} and
\texttt{three\_step\_flow\_jacobian} by the two cases $d=2,3$.
The map is polynomial, $C^1$, and has nonzero determinant on the open set
$\Omega$. The injective change-of-variables inequality gives
\[
 |E_\ell|\geq\int_{\mathcal P}|\det D\Phi|
 \geq 2^{-15}\alpha_1^2\alpha_d^2.                 \tag{J2}
\]
One may apply the change-of-variables theorem with the target integrand
$\mathbf1_{E_\ell}$; no separate measurable-image lemma for
$\Phi(\mathcal P)$ is needed. Substituting $\alpha_i=K/|E_i|$ gives
\[
 K^4\leq2^{15}|E_1|^2|E_d|^2|E_\ell|.
\]
Taking fourth roots and using $2^{15/4}\leq16$ proves (R).
\end{proof}

\begin{lemma}[Strong $L^1$ bound from the restricted three-set bound]
\label{lem:three-set-strong-upper-half}
Fix a distinguished input $j$. Choose
\[
 (d,a,\ell)=
 \begin{cases}
  (2,2,3),&j=1,\\
  (2,1,3),&j=2,\\
  (3,1,2),&j=3.
 \end{cases}
\]
Thus $\{j,a\}=\{1,d\}$. There is a constant $C_1>0$ such that, for all measurable inputs with finite displayed norms,
\[
 \|T(f_1,f_2,f_3)\|_1
 \leq C_1\|f_j\|_2\|f_a\|_{8/3}\|f_\ell\|_4.       \tag{S}
\]
\end{lemma}

\begin{proof}
It is enough to consider nonnegative functions. Let
$\Lambda(f)=\int T(f)$ and set
\[
 v_1=v_d=\tfrac12,\quad v_\ell=\tfrac14,
 \qquad b_j=\tfrac12,\quad b_a=\tfrac38,\quad b_\ell=\tfrac14,
 \qquad p_i=1/b_i.
\]
The preceding lemma gives the restricted form bound at $v$. The elementary
bounds $\Lambda(\mathbf1_E)\leq |E_i|$ give the bounds at the three standard
vectors $e_i$. The target lies strictly inside their tetrahedron:
\[
 b=\tfrac12v+\tfrac14 e_j+\tfrac18e_a+\tfrac18e_\ell.
\]
Here is a direct finite-level summation proof, so this step need not rely
on a general restricted-weak-type interpolation statement.

Put $\eta=1/64$. For every $\sigma\in\{-1,1\}^3$, $c=b+\eta\sigma$ is in
the same tetrahedron. Explicit barycentric coefficients are
\[
 \tau=4\Bigl(\sum_i c_i-1\Bigr),\qquad
 \lambda_i=c_i-\tau v_i.
\]
They satisfy $\tau>0$, $\lambda_i>0$, and $\tau+\sum_i\lambda_i=1$.
For example, uniformly over the signs and the three choices of $j$,
$\tau\geq5/16$ and $\lambda_i\geq3/64$. These are eight rational checks
for each $j$. Taking the weighted geometric mean of the four restricted
bounds proves
\[
 \Lambda(\mathbf1_{E_1},\mathbf1_{E_2},\mathbf1_{E_3})
 \leq16\prod_i|E_i|^{b_i+\eta\sigma_i}.             \tag{S1}
\]
Zero-measure sets are handled before any division or logarithm.

For finite-valued, finite-measure-support, nonnegative inputs with
$\|f_i\|_{p_i}=1$, use the value-level sets
\[
 E_{i,n}=\{x:2^n\leq f_i(x)<2^{n+1}\},\qquad n\in\mathbb Z.
\]
Only finitely many such sets are nonempty and
$|E_{i,n}|\leq2^{-p_i n}$. For each triple $(n_1,n_2,n_3)$ choose
$\sigma_i=1$ if $n_i\geq0$, and $\sigma_i=-1$ otherwise.
Since all exponents in (S1) are positive,
\[
 2^{n_1+n_2+n_3}\Lambda(\mathbf1_{E_{1,n_1}},
     \mathbf1_{E_{2,n_2}},\mathbf1_{E_{3,n_3}})
 \leq16\prod_i2^{-\eta p_i|n_i|}.                  \tag{S2}
\]
Dominate $f_i$ by $\sum_n2^{n+1}\mathbf1_{E_{i,n}}$, expand the positive
trilinear form, and sum (S2). A possible finite constant is
\[
 C_1=128\prod_{i=1}^3
       \frac{1+2^{-\eta p_i}}{1-2^{-\eta p_i}}.
\]
Rescale to remove normalization; a zero input norm is handled separately.
Monotone simple approximation and spatial truncation extend (S) to all
nonnegative inputs. For this passage define the positive average by nonnegative
integration, and convert to real-valued functions only after proving finiteness.
For complex inputs use $|T(f)|\leq T(|f|)$.
This proof uses dyadic \emph{values}, not dyadic spatial resolution of sets.
The proposed Lean declaration is
\texttt{restricted\_form\_interior\_strong}.
\end{proof}

\begin{lemma}[Subunit bound for the actual monomial system]
\label{lem:adjoint-gain-to-subunit}
For each $j\in\{1,2,3\}$ and all measurable inputs with finite displayed norms, the upper-half monomial average at scale one
satisfies
\[
 \|\widetilde A_1(f_1,f_2,f_3)\|_{8/9}
 \leq C_j\prod_{i=1}^3\|f_i\|_{p_i},               \tag{Q}
\]
where
\[
 (p_1,p_2,p_3)=
 \begin{cases}
 (2,8/3,4),&j=1,\\
 (8/3,2,4),&j=2,\\
 (8/3,4,2),&j=3.
 \end{cases}
\]
In particular $p_j=2$, every $p_i>1$, and
$\sum_i1/p_i=9/8=1/q$ for $q=8/9\in(1/2,1)$.
The historical label is retained only for compatibility: no implication
for arbitrary polynomial translation averages is asserted.
\end{lemma}

\begin{proof}
Let $q=8/9$ and use the strong $L^1$ bound (S). Partition $\mathbb R^3$
into $Q_z=z+[0,1)^3$, $z\in\mathbb Z^3$, and set
$R_z=z+[0,2)^3$ and $f_{i,z}=f_i\mathbf1_{R_z}$. For $x\in Q_z$ and
$t\in I$, every $x+t^i e_i$ lies in $R_z$, so
\[
 \mathbf1_{Q_z}T(f)=\mathbf1_{Q_z}T(f_{1,z},f_{2,z},f_{3,z}).
\]
For $g\geq0$ and $|Q_z|=1$, ordinary H\"older (or concavity of the
$q$-th power on a probability space) gives
$\int_{Q_z}g^q\leq(\int_{Q_z}g)^q$.
With $M_{i,z}=\int_{R_z}|f_i|^{p_i}$, (S) therefore implies
\[
 \int_{Q_z}|T(f)|^q\leq C_1^q\prod_iM_{i,z}^{q/p_i}.
\]
Sum over $z$. Discrete H\"older has exponents $p_i/q>1$, whose reciprocals
sum to one. Also every point belongs to at most eight of the sets $R_z$.
Consequently
\begin{align*}
 \int_{\mathbb R^3}|T(f)|^q
 &\leq C_1^q\sum_z\prod_iM_{i,z}^{q/p_i}\\
 &\leq C_1^q\prod_i\Bigl(\sum_zM_{i,z}\Bigr)^{q/p_i}\\
 &\leq 8C_1^q\prod_i\|f_i\|_{p_i}^{q}.
\end{align*}
Thus $C_j=8^{1/q}C_1$ is admissible. In Lean, prove the discrete H\"older
step first for finite sets of $z$, and then pass to the countable sum by
monotone convergence. Proposed declarations are
\texttt{local\_L1\_to\_balanced\_subunit} and
\texttt{monomial\_subunit\_upperHalf}.
There is no absorption remainder and no use of duality for $L^q$.
\end{proof}

\begin{corollary}[Internal subunit endpoint]
\label{cor:kosz-subunit-internal}
For each distinguished input $j$, the subunit endpoint required in the
blueprint holds for the full monomial averages, without assuming
(6.24) or deducing a quasi-Banach estimate from an unspecified adjoint gain.
One can take $B=9/8$ and $(b_i)=(1/p_i)$ as in (Q).
\end{corollary}

\begin{proof}
Apply (Q). The anisotropic scaling and the upper-half to full-average
dyadic summation already given in the proof of
\texttt{thm:kosz-subunit} apply unchanged: $\sum_i1/p_i=1/q$ cancels the
scaling powers, and
\[
 \sum_{n\geq0}2^{-(n+1)q}=\frac1{2^q-1}<\infty.
\]
Thus the full-average constant can be taken as
$C_j(2^q-1)^{-1/q}$ and then enlarged to a dyadic constant as before.
The proof dependency is now: three measurable refinements, fibre
separation, the explicit injective flow, change of variables, finite
value-level summation, and fixed-scale cube localization. In particular,
Corollary~5.3 is not used as an abstract bridge across $q=1$.
\end{proof}

\subsection*{Why the abstract implication must be removed}
The obstruction is to the generic lemma, not to the monomial estimate just
proved. Consider the equal-linear corner averages
\[
 B_N(f)(x)=N^{-1}\int_0^N\prod_{i=1}^3 f_i(x+t e_i)\,dt.
\]
They meet the stated polynomial-degree restriction. For their positive
$j$-adjoint $S_N$, one has
\[
 \|S_N(g_0,g_a,g_b)\|_2\leq N^{-3/4}
        \|g_0\|_4\|g_a\|_4\|g_b\|_4.             \tag{C1}
\]
To verify this, the four-dimensional Loomis--Whitney inequality, which
follows by repeated H\"older, gives
\[
 \int_{\mathbb R^4}h_0(x)\prod_{i=1}^3 h_i(x+t e_i)\,dx\,dt
 \leq\prod_{i=0}^3\|h_i\|_3.
\]
An explicit unit-determinant reduction to coordinate projections is
$x=(-z_1,-z_2,-z_3)$, $t=z_0+z_1+z_2+z_3$; each of the four factors then
omits its corresponding $z_i$, and its remaining three-dimensional linear
change of variables also has absolute determinant one. Restricting $t$
and dualizing gives
$\|S_N(g)\|_{3/2}\leq N^{-1}\prod_i\|g_i\|_3$.
Parameter H\"older on the probability interval gives
\[
 S_N(g)\leq S_N(g_0^{4/3},g_a^{4/3},g_b^{4/3})^{3/4},
\]
which proves (C1). Thus all three adjoint input exponents are $4$, their
reciprocal sum is $3/4=1/(4/3)$, and the scale gain is strictly positive.

Nevertheless, no balanced subunit estimate can hold for $B_1$. Put
\[
 E_\delta=\{x\in[-2,2]^3:0\leq x_1+x_2+x_3\leq\delta\},
 \quad f_i=\mathbf1_{E_\delta},\quad 0<\delta<1/8.
\]
Then $|E_\delta|\leq16\delta$. On the fixed cube
$U=[-1/4,-1/8]^3$, the common condition for the three translated indicators
is $0\leq x_1+x_2+x_3+t\leq\delta$; the corresponding interval lies in
$(0,1)$ and has length $\delta$. Hence $B_1(f)=\delta$ on $U$.
An estimate with $B=\sum_i1/q_i=1/q>1$ would force
\[
 \delta |U|^{1/q}\leq C(16\delta)^B
\]
for every small $\delta$, which is impossible. This explains why a correct
repair uses the actual monomial geometry rather than filling in an
abstract recursion from the adjoint gain alone.

\subsection*{Source and formalization boundary}
Source: D.~Kosz, M.~Mirek, S.~Peluse, R.~Wan, J.~Wright,
\emph{The multilinear circle method and a question of Bergelson},
arXiv:2411.09478v4, supplied PDF. Relevant locations: pp.~68--69
(measurable refinements and Lemma~5.14), pp.~69--70 (parameter towers),
pp.~72--74 (real change of variables and separation), and p.~79,
equation~(6.24).

The analytic dependencies still to implement in Lean are ordinary Tonelli
and translation invariance, measurable parameter integration, injective
$C^1$ change of variables, monotone approximation, H\"older, and geometric
series. The new flow's injectivity and Jacobian, the three updates, the
barycentric coefficients, and the final summation are specified above;
none is an unnamed stopping-time or quasi-Banach interpolation oracle.
\end{document}
